\begin{document}

\section{提案手法}

\subsection{パイプライン}

図\ref{figure:pipeline}に，本研究で用いるCLOUSEAU型のエージェント型ログ探索パイプラインを示す．
入力は，CBC alert rows，またはhost/process/time windowである．
CLOUSEAU本来のPOI起点調査\cite{clouseau}にならい，Chief agentは初期手がかりから調査方針を生成する．
Investigation agentは，確認すべき事実をQAAgent / SQL expertへの質問に変換する．
QAAgent / SQL expertは，SQLite化されたエンドポイントログに問い合わせ，timestamp，source stream，PID/PPID，process tree，command line，registry path，source row idなどを返す．
Final synthesisは観測結果を統合し，主行動列，補助証跡，limitationsを出力する．

\begin{figure}[ht]
    \centering
    \fbox{
    \begin{minipage}{0.92\linewidth}
        \small
        \centering
        SOC clue / CBC alert / process-time clue\\
        $\downarrow$\\
        Chief agent: 調査方針を生成\\
        $\downarrow$\\
        Investigation agent: 確認事項を質問へ変換\\
        $\downarrow$\\
        QAAgent / SQL expert: 観測行を取得\\
        $\downarrow$\\
        Final synthesis: 主行動列と近傍行動を分離\\
        $\downarrow$\\
        行動ステップ列 + evidence + limitations
    \end{minipage}}
    \caption{SOC調査起点から証跡付き行動列を復元するパイプライン}
    \label{figure:pipeline}
\end{figure}

本手法の設計上の要点は，仮説生成と観測ログ検索を分離することである．
Chief agentは直接SQLを書かず，Investigation agentが観測済みの値に基づいて次の確認事項を作る．
DBに問い合わせるのはQAAgent / SQL expertであり，Final synthesisは観測行で支えられる行動だけを主系列に含める．
これにより，alert nameやreasonの言い換えを行動復元と誤認しないようにする．

\subsection{出力と評価単位}

出力は，\texttt{steps}，\texttt{sequence}，\texttt{evidence}，\texttt{nearby}，\texttt{limitations}からなる．
\texttt{steps}は主体，操作，対象，コマンドライン，証跡を持つ主行動ステップである．
\texttt{sequence}は観測された操作を短く並べた系列であり，アラート名ではなくログ上の操作を表す．
\texttt{evidence}はsource stream，row id，timestamp，PID，command lineなどの根拠である．
\texttt{nearby}は時刻や対象が近いが主系列に含めない行動を明示する．

採点では，gold chainを行動itemに分解する．
必須itemは，ログから直接確認できる主体，操作，対象，コマンドライン，critical evidenceである．
観測行だけから主張できない悪性判定，ユーザ意図，業務目的，ファイル内容の推定は必須itemに含めない．
候補出力の主張は正規化され，gold itemと照合される．
主指標は，行動item recall，critical evidence recall，sequence order，candidate claim precision，overclaim slot countである．

\end{document}
